\scanpage{262}
\[
\tag{93}
\diagram{\begin{tikzcd}[column sep=huge,row sep=huge]
\underline{\mathrm{germesdiff}}(X,x)
\arrow[r,"{\substack{\gamma\longrightarrow\underline\Delta(\gamma)\\\approx}}"]
\arrow[d,"\approx","{\substack{\gamma\mapsto\mathrm{Revuni}(D_x^*,\gamma)\\=D_{x,\gamma}^{*\sim}}}"']
& \Pi_1(\mathbb S_x)
\arrow[d,"{\approx\quad a\mapsto\mathrm{Revuni}(\mathbb S_x,\cdot)}"]
\arrow[dl,dashed,"{\Theta\ \approx}"'] \\
\underline{\mathrm{Rev.uni.}}(D_x^*) \arrow[r,"\approx"]
& \mathrm{Revuni.}(\mathbb S_x)
\end{tikzcd}}
\]
\margin{Of course, the difficulty here came from the fact that we had no
direct morphism between the germ $D_x^*$, and the space $\mathbb S_x$...}

\unsure{The} only functor that has not yet made explicit is the functor
\[
\tag{94}
\diagram{\begin{tikzcd}[column sep=small,row sep=small]
\underline{\mathrm{Revuni}}(D_x^*)
\arrow[rr,"\approx"] \arrow[dr,dashed,"{\Theta^{-1}}"']
& & \mathrm{Revuni}(\mathbb S_x) \\
& \Pi_1(\mathbb S_x) \arrow[ur,"{\mathrm{can}}"'] &
\end{tikzcd}}
\]
defined as the composite of $\Theta^{-1}$ with the canonical functor
$\Pi_1(\mathbb S_x)\longrightarrow\mathrm{Rev.uni.}(\mathbb S_x)$.

Finally, in these considerations of differential topological geometry,
the key construction is that of the functor $\Theta$ (81).
We have taken some pains to define it in an invariant way.
In the case where we have a local isomorphism
\[
\tag{90}
f:\bigl(\underbrace{\mathrm{Tang}_x(X)}_{V},0\bigr)\longrightarrow(X,x),
\]
the construction of \unsure{this|---} $\Theta$ can be made explicit as
a canonical construction in a vector space $V$.
Every homogeneous half-line

\scanpage{263}
\mspage{779}

$\Delta$ in $V$ defines a germ of arcs (\unsure{differentiable}) in $V$ issuing from $0$, whence a corresponding universal covering $D_{V,0;\Delta}^{*\sim}$ of $D_{V,0}^{*}$, hence a universal covering $D_{x,\Delta}^{*\sim}$ of $D_x^*$.
We must see how this $D_{x,\Delta}^{*\sim}$ depends ``continuously'' on $\Delta$, and what amounts to the same thing, how the universal covering $V_{\Delta}^{*\sim}$ of $V^*=V\setminus\{0\}$ depends continuously on $\Delta$, which is quite \unsure{standard}\unsure{...}.

Now let \inserted{$\unsure{u|\varphi}$} be \unsure{an} an \deleted{diffeomorphism}\inserted{\unsure{automorphism}} of $X$, fixing $x$, and differentiable at $x$. We thus have \deleted{\unsure{some}} \unsure{automorphisms} induced \deleted{of}
\[
\tag{\unsure{95|94}}
{}^{\unsure{u|\varphi}}D_x^*\quad\text{automorphism of }D_x^*,\qquad
{}^{\unsure{u|\varphi}}\mathbb S_x\quad\text{automorphism of }\mathbb S_x.
\]

Consider on the other hand a universal covering $\widetilde{\mathbb S}_x$ of $\mathbb S_x$, and let $D_x^{*\sim}$ be an associated universal covering of $D_x^*$, thanks to the equivalence (94). I say that we then have a

\scanpage{264}

canonical bijection
\[
\tag{96}
\diagram{\begin{array}{c@{\qquad}c@{\qquad}c}
\text{set of liftings}&\simeq&\text{set of liftings}\\
\text{of }{}^uD_x^*\text{ to an}&&\text{of }{}^u\mathbb S_x\text{ to an}\\
\text{automorphism }{}^uD_x^{*\sim}\text{ of the}&&\text{automorphism }{}^u\mathbb S_x^\sim\text{ of}\\
\text{germ }D_x^{*\sim}&&\mathbb S_x^\sim
\end{array}},
\]
\margin{the construction (96) will be carried out in the following pages, leading to (122) page \unsure{783}}
bijection which is an isomorphism of $\overline T_x$-torsors.
(where $\overline T_x$ is the $\mathbb Z$-orientation module of $X$ at $x$, let us recall).

\margin{E) \uline{Axiomatization} : \unsure{the structure|structure} of \uline{correlation} between universal coverings $\mathbb S_x^\sim$ \uline{and} $D_x^{*\sim}$.}
I feel like giving an axiomatic description of (94) and (96).

1) Let \inserted{$D_x=$} $(X,x)$ be a germ of a differentiable surface, whence
$V_{X,x}=V_x=\mathbb V=\operatorname{Tang}_x(X)$,
and $\mathbb S_x=V^*/\mathbb R^{*+}$ (where $\mathbb V^*=V\setminus\{0\}$), as well as $D_{X,x}^*=D_x^*$.
\margin{``germ'' : we \unsure{tacitly include} the filter of neighborhoods of $x$ in $X$}
\deleted{Let $\mathbb S_x^\sim$ (\unsure{resp.} $D_x^{*\sim}$) be \unsure{universal} coverings of $\mathbb S_x$ (resp. $D_x^*$).}
\inserted{We denote by}
\deleted{$\unsure{\mathcal A|\in}$}
$\mathcal A_{X,x}$ \unsure{or} $\mathcal A_x$ the set of \uline{arcs} of $\mathbb S_x$, in one-to-one correspondence with the set $\mathcal C_{X,x}=\mathcal C_x$ of
\deleted{``\unsure{angular sectors}'' of $V$ \unsure{proper}}
\inserted{``closed cones'' \unsure{proper}} of $V$ ($\ne V$ and $\{0\}$).
For $A\in\mathcal A_x$ or $C\in\mathcal C_x$, we denote by
\deleted{$\unsure{\widetilde A|f(A)}$ \unsure{or} $\unsure{\widetilde C|f(C)}$}
\inserted{$\unsure{\mathrm{Ep}|E_p}(A)$ or $\unsure{\mathrm{Ep}|E_p}(C)$}
the corresponding germ of spaces over $D_x^*$.
\margin{``\unsure{thickening}'' \deleted{\unsure{image|germ}} of $A$ (or of $C$) in $D_x^*$}
If $\widetilde{\mathbb S}$ is a universal covering of $\mathbb S$, $\widetilde D_x^*$ a universal covering of $\widetilde D_x^*$, we consider the sets

\scanpage{265}
\mspage{780}

\[
\tag{97}
\diagram{\begin{array}{r@{\,}c@{\qquad}c@{\qquad}r@{\,}c}
&\Gamma\bigl(\deleted{\unsure{\mathbb S|\widetilde A}}A,\mathbb S_x^\sim/\mathbb S_x\bigr)&,&
&\Gamma\bigl(\inserted{\unsure{\mathrm{Ep}|E_p}}(\widetilde A),D_x^{*\sim}/D_x^*\bigr)\\
&\Vert&&&\Vert\\
&\operatorname{Hom}_{\mathbb S}(A,\mathbb S_x^\sim)&&&\operatorname{Hom}_{D_x^*}(\widetilde A,D_x^{*\sim})\\
\mathrm{déf}&\Vert&&\mathrm{déf}&\Vert\\
&\mathbb S_x^\sim(A)&&&D_x^{*\sim}\bigl(\unsure{\mathrm{Ep}|E_p}(A)\bigr)
\end{array}}
\]
which are both Torsors under the same group $\overline T_{X,x}=\overline T_x$ (the $\mathbb Z$-module of orientations of $X$ at the point $x$). Consider the set
\[
\tag{98}
H(A)=\operatorname{Isom}_{\unsure{T_x|\overline T_x}}\bigl(\mathbb S_x^\sim(A),
\deleted{\unsure{D_x^{*\sim}(\widetilde A)}}\unsure{)}
\inserted{D_x^{*\sim}\bigl(\unsure{\mathrm{Ep}|E_p}(A)\bigr)}\bigr)
\]
which is itself a $\overline T_x$-torsor. If we have \deleted{\illegible{1}} an inclusion
\[
A\subset A'
\]
we deduce from it, by restriction, isomorphisms of $\overline T_x$-torsors
\[
\tag{99}
\mathbb S_x^\sim(A')\xrightarrow{\sim}\mathbb S_x^\sim(A),\qquad
D_x^{*\sim}\bigl(\unsure{\mathrm{Ep}|E_p}(A')\bigr)
\unsure{-|\longrightarrow}
D_x^{*\sim}\bigl(\unsure{\mathrm{Ep}|E_p}(A)\bigr)
\]
whence a canonical isomorphism
\[
\tag{100} H(A)\xrightarrow{\sim}H(A')
\]
so that \deleted{$A\unsure{\longmapsto|\longleftarrow}H(A)$ \unsure{defines}}
\[
\tag{101}
A\longmapsto H(A),\quad\text{functor}\quad
\underline{\mathcal A}_x\xrightarrow{H}\underline{\operatorname{Tors}}(\overline T_x),
\]
where $\underline{\mathcal A}_x$ is the ordered category associated to the ordered set $\mathcal A_x$. We thus find by factorization a functor \unsure{from the|from} \deleted{the category} enveloping groupoid of $\underline{\mathcal A}_x$ to $\underline{\operatorname{Tors}}(\overline T_x)$;
\[
\tag{102}
\diagram{\begin{tikzcd}[column sep=large,row sep=large]
&\underline{\mathcal A}_x\arrow[dl,"{\unsure{\rho|\gamma}}"']\arrow[dr,"H"]&\\
\underline{\operatorname{Gr}}(\underline{\mathcal A}_x)\arrow[rr,"{\underline{\operatorname{Gr}}(H)}"]&&\underline{\operatorname{Tors}}(\overline T_x)
\end{tikzcd}}.
\]

\scanpage{266}

\deleted{I say that}\inserted{We know that} $\underline{\operatorname{Gr}}(\underline{\mathcal A}_x)$ is \deleted{\unsure{isomorphic}} a $\overline T_x$-groupoid, more precisely the functor
\[
\deleted{\operatorname{Hom}_{\underline{\operatorname{Gr}}(\underline{\mathcal A}_x)}(A,A')\xrightarrow{\sim}}
\]
\[
\tag{103}
\diagram{$\left\{\begin{array}{rcl}
\underline{\mathcal A}_x&\longrightarrow&\underline{\operatorname{Revuni}}(\mathbb S_x)\overset{\mathrm{déf}}{=}\mathbb S_{x,A}^{\deleted{\unsure{*}}\sim}\\[1ex]
A&\longmapsto&\deleted{\unsure{\mathbb S}}\operatorname{Revuni}(\mathbb S,A)
\end{array}\right.$}
\]
which factors as a \deleted{functor}\inserted{morphism} of groupoids
\[
\tag{104}
\underline{\operatorname{Gr}}(\underline{\mathcal A}_x)\xrightarrow{(\simeq)}\underline{\operatorname{Rev.uni.}}(\mathbb S_x)
\]
gives an \uline{equivalence of categories} (104), hence
\[
\tag{105}
\operatorname{Hom}_{\underline{\operatorname{Gr}}(\underline{\mathcal A}_x)}(A,A')\xrightarrow{\sim}
\operatorname{Hom}_{\deleted{\unsure{\underline{\operatorname{Revuni}}}}\inserted{\unsure{\mathbb S_x}}}(\mathbb S_{x,A}^{\sim},\mathbb S_{x,A'}^{\sim})
\]
\margin{\unsure{i.e. we have for} $A,A'\in\mathcal A_x$, \unsure{the canonical isomorphism}
\[
\operatorname{Hom}_{\underline{\operatorname{Gr}}(\underline{\mathcal A}_x)}(A,A')\simeq
\operatorname{Isom}_{\underline{\operatorname{Revuni}}(\mathbb S_x)}(\mathbb S_{x,A}^{\sim},\mathbb S_{x,A'}^{\sim}).
\]}

This \deleted{\unsure{done|being}} recalled, I say that the functor $\underline{\operatorname{Gr}}(H)$ of (102) is \uline{essentially} \unsure{constant} \inserted{i.e. it defines on the $\pi_1$ the zero homomorphism $\overline T_x\xrightarrow{0}\overline T_x$} --- or what amounts to the same thing, there is between the $H(A)$, \inserted{including the isomorphisms (100),} a \uline{transitive system} of \uline{isomorphisms}. We then denote \deleted{by} the ``common value'' of these $H(A)$, which depends only on the choice of $\mathbb S_x^{\sim}$ and of $D_x^{*\sim}$, by $H(\mathbb S_x^{\sim},D_x^{*\sim})$. \deleted{\unsure{We thus have}}\inserted{This set} is endowed with the structure defined by isomorphisms of $\overline T_x$-\unsure{torsors|functors}
\[
\tag{106}
\unsure{i_A|\widetilde{\imath}_A}:H(A)\xrightarrow{\sim}H(\mathbb S_x^{\sim},D_x^{*\sim})
\]
satisfying the compatibility condition, express-